\section{Method}

The method models reported results with a batch/scanner grouping factor. The
response can be a vote share, a transformed vote share, or a residual from a
baseline model that already accounts for contest, jurisdiction, ballot style,
mode, time, or comparable historical information. The grouping factor is the
batch, scanner, drop box, upload, or other operational unit being evaluated.

A fixed-effect version estimates a separate coefficient for each group. A
random-effect version treats group effects as draws from a common distribution
and estimates the degree to which each group departs from the population of
comparable groups. Either form asks the same forensic question: is the estimated
group effect large relative to the variation observed within and across groups?
A report may rank groups by standardized effect size, posterior or empirical
shrinkage estimate, residual dispersion, or a calibrated tail probability.

The score is not a fraud finding. Legitimate operational differences routinely
produce real batch effects. Vote-by-mail batches can have different candidate
shares than in-person precinct batches. Early batches can differ from
election-day batches. A scanner assigned to one geography or ballot style can
inherit that geography or style. For that reason, the preferred baseline makes
the comparison population explicit and includes the operational variables needed
to avoid treating ordinary election administration as an anomaly.

The method flags groups whose effect remains large after that comparison. It
also records the scale on which the score was computed, the baseline population,
and an investigation note describing the most likely benign grouping explanation
when one is visible from the inputs. The statistical machinery is established;
this paper fixes the W-series framing and report shape for this analytic family.
